% =====================================================================
% 5. Conclusão
% Sem informação nova e sem citação.
% =====================================================================
\section{Conclusão}\label{sec:conclusao}

Neste trabalho, o grupo fez todos os passos para desenhar uma cena 3D na tela. Foram criados cinco
objetos do dia a dia, e cada um virou uma malha de triângulos. Depois, os objetos foram colocados
juntos na cena, vistos por uma câmera, projetados em perspectiva, iluminados e pintados em três
tamanhos de imagem.

Os resultados mostram o que cada passo faz. Mudar o \textit{eye} muda o ponto de vista. A luz dá a
sensação de volume aos objetos. O Z-buffer faz o objeto da frente esconder o de trás. Quanto maior a
imagem, mais lisos ficam os contornos.

Ainda há o que melhorar. Os objetos não foram girados nem aumentados na cena, a lixeira aparece
cortada na câmera principal e o \textit{eye} só muda no código. \confirmar{ajustar esta lista conforme o código final}
Como próximos passos, o usuário poderia escolher o \textit{eye} sem mexer no código, e os polígonos
fora do volume de visão poderiam ser cortados.
